\subsection{The all-weight gate at every regular coefficient prime}

\begin{result}{E.209}{Proved}
{Regular coefficients have no fixed-ramification second obstruction in any weight}
Let $\ell$ be an odd prime, let $K=\mathbb Q(\zeta_\ell)$, and
assume $\ell\nmid h_K$, where $h_K$ is the ideal class number.
Thus ``regular'' in this statement means precisely this class-number
condition; no Bernoulli-number criterion is needed. Put
$S=\{\ell,\infty\}$ and $G_S=G_{\mathbb Q,S}$.
For every integer $w$ and every $n\ge1$,
\[
 H^2(G_S,\mathbb Z/\ell^n(w))=0,
 \qquad H^2_{\rm cts}(G_S,\mathbb Z_\ell(w))=0.              \tag{E.209.1}
\]
In particular these conclusions hold for $\ell=5$, because the
class-number-one calculation for $\mathbb Q(\zeta_5)$ was given
in E.202. For $w\ge2$, the strict localization map of A.22 is
therefore an isomorphism
\[
 H^2_{\rm cts}(G_{\mathbb Q},\mathbb Z_\ell(w))
 \simeq
 \frac{c_0(\{p:p\ne\ell\},\mathbb Z_\ell)}
 {\{((p^{w-1}-1)y_p)_p:(y_p)_p\in c_0(\mathbb Z_\ell)\}}.
                                                               \tag{E.209.2}
\]
The denominator is the actual image, without closure.

For the actual pro-$\ell$ geometric quotients of the tripod, every
class-$r$ section, $r\ge2$, whose abelian pair is ordinary and
which matches its ordinary point section at each support prime
$p\ne\ell$ has a global continuous integral class-$(r+1)$ lift.
Preserving all these matches is exactly the separate localization
condition of B.26. Thus that condition is the sole remaining
inductive obstruction at every depth under the displayed
class-number hypothesis. Existence of a compatible infinite
sequence of matching lifts is not asserted.
\end{result}

\begin{proof}
We first prove the finite assertion over $K$, including the
inflation comparison. Let $S_K$ contain its unique prime over
$\ell$ and its complex places, and write
\[
 1\longrightarrow N\longrightarrow G_K
 \longrightarrow G_{K,S_K}\longrightarrow1.
\]
The coefficient $\mu_\ell$ is a trivial $G_K$-module.
The five-term cohomology sequence contains
\[
 \begin{gathered}
 H^1(G_K,\mu_\ell)\longrightarrow
 H^1(N,\mu_\ell)^{G_{K,S_K}}
 \longrightarrow H^2(G_{K,S_K},\mu_\ell)\\
 \xrightarrow{\rm inf} H^2(G_K,\mu_\ell).
 \end{gathered}                                                 \tag{E.209.3}
\]
We prove that its first arrow is surjective rather than simply
assuming injectivity of its last arrow.

An invariant continuous character $\varphi:N\to\mu_\ell$
has finite-support tame residue data outside $S_K$.
Indeed continuity supplies an open normal subgroup $U$ of
$G_K$ with $N\cap U\subseteq\ker\varphi$. The finite
extension fixed by $U$ ramifies at only finitely many primes.
At every other prime outside $S_K$ its full inertia is in
$N\cap U$, and the residue of $\varphi$ is zero.
At a prime $\mathfrak p\nmid\ell$, a homomorphism from inertia
to $\mu_\ell$ kills wild inertia and factors through the
cyclic tame $\ell$-quotient. Fix its generator so that the
Kummer character of a uniformizer has value one in
$\mathbb F_\ell$. Thus $\varphi$ determines a finite tuple
$(v_{\mathfrak p})_{\mathfrak p\notin S_K}$.
These residues determine $\varphi$: by the definition of the
maximal extension unramified outside $S_K$, $N$ is the closed
normal subgroup generated by these inertia groups, and
$\varphi$ is invariant under conjugation by $G_K$.

Every such finite tuple is the residue tuple of a global
Kummer class. Choose integral representatives of its entries
and form the fractional ideal
$I=\prod_{\mathfrak p\notin S_K}\mathfrak p^{v_{\mathfrak p}}$.
The ideal $I^{h_K}$ is principal, say $(\beta)$.
Choose an integer $a$ with $a h_K\equiv1\pmod\ell$.
Then $\alpha=\beta^a\in K^\times$ has exactly the prescribed
valuations modulo $\ell$ at all primes outside $S_K$.
Kummer theory gives $\kappa_\ell(\alpha)\in H^1(G_K,\mu_\ell)$;
its restriction to $N$ agrees with $\varphi$ on every inertia
group and hence on $N$. This proves surjectivity in E.209.3.

Consequently $H^2(G_{K,S_K},\mu_\ell)$ injects into
$H^2(G_K,\mu_\ell)=\operatorname{Br}(K)[\ell]$.
At every finite place outside $S_K$ an inflated class factors
through the unramified quotient. Its second cohomology is
zero, so its local Brauer invariant is zero. The complex
places contribute nothing. There is only one remaining
finite place, the unique place above $\ell$. The global
Brauer injection and the sum-of-invariants relation force
that invariant to vanish too
\citep[II.2.1--2.2, pp.~163--165]{CohMilneADT}.
It follows that
\[
                  H^2(G_{K,S_K},\mu_\ell)=0.                    \tag{E.209.4}
\]
This argument is a literal Galois-cohomology inflation
argument and does not require an unproved comparison with
the cohomology of an arithmetic scheme.

Every finite extension of $K$ unramified outside $S_K$ has
normal closure over $\mathbb Q$ unramified outside $S$:
the conjugate extensions and their compositum remain
unramified away from $\ell$. Hence the corresponding
maximal extensions give the exact sequence
\[
 1\longrightarrow G_{K,S_K}\longrightarrow G_S
 \longrightarrow\Delta\longrightarrow1,
 \qquad\Delta=\operatorname{Gal}(K/\mathbb Q).
\]
The order of $\Delta$ is $\ell-1$, invertible on any
$\ell$-primary module. Its higher cohomology vanishes by
averaging. Every $\mathbb F_\ell(w)$ is trivial on
$G_{K,S_K}$, so E.209.4 and restriction give
\[
 H^2(G_S,\mathbb F_\ell(w))
  =H^2(G_{K,S_K},\mathbb F_\ell(w))^\Delta=0
 \quad(w\in\mathbb Z).                                        \tag{E.209.5}
\]

Fix $T=\mathbb Z_\ell(w)$ and write $T_n=T/\ell^nT$.
The actual equivariant short exact sequence
\[
 0\longrightarrow T_1\xrightarrow{\ \ell^n\ }T_{n+1}
 \longrightarrow T_n\longrightarrow0
\]
has first map $t\bmod\ell T\mapsto\ell^nt\bmod\ell^{n+1}T$.
Its degree-two exact segment and E.209.5 prove
$H^2(G_S,T_n)=0$ by induction on $n$.
For each $n$, $H^1(G_S,T_n)$ is finite: its cocycles give
finite semidirect-product extensions of bounded degree
ramified only in $S$, and fixed-ramification finiteness
applies \citep[Theorem8.42, p.~151]{CohMilneANT}.
Thus their inverse system is Mittag--Leffler. The actual
continuous-cohomology sequence
\citep[(2.1) and Theorem2.2, p.~215]{CohJannsen1988}
\[
 0\longrightarrow\varprojlim{}^1 H^1(G_S,T_n)
 \longrightarrow H^2_{\rm cts}(G_S,T)
 \longrightarrow\varprojlim H^2(G_S,T_n)\longrightarrow0
\]
has both outside terms zero. This proves the integral
assertion separately from its finite reductions.

Equation E.209.2 follows from the exact strict localization
sequence of A.22 with its first term zero. For the tower
claim, B.25 puts the next obstruction of the specified
boundary-matched class-$r$ section in
$H^2_{\rm cts}(G_S,\mathbb Z_\ell(r+1))^{b_{r+1}}$.
It vanishes by E.209.1. B.26 gives the exact additional
condition for a choice of lift to retain every match.
We have removed the fixed-ramification obstruction, not
that simultaneous matching condition or the absolute
strict-localization group.
\end{proof}

\begin{result}{E.210}{Proved}
{The exact global correction parameters and coefficient-place matching at regular primes}
Retain the hypotheses of E.209 and let $w\ge2$. Inflation
identifies the global integral first cohomology with its
fixed-ramification version, and
\[
 \begin{aligned}
 H^1_{\rm cts}(G_{\mathbb Q},\mathbb Z_\ell(w))
 &=H^1_{\rm cts}(G_S,\mathbb Z_\ell(w))\\
 &\simeq
 \begin{cases}
  \mathbb Z_\ell,&w\text{ odd},\\
  \mathbb Z/\ell^{1+v_\ell(w)},&\ell-1\mid w,\\
  0,&w\text{ even and }\ell-1\nmid w.
 \end{cases}
 \end{aligned}                                                \tag{E.210.1}
\]
In the middle case a generator is
\[
              \eta_w(g)=\frac{\chi_\ell(g)^w-1}
                               {\ell^{1+v_\ell(w)}}.           \tag{E.210.2}
\]
At the coefficient place, localization is an isomorphism
onto the entire $H^1(G_{\mathbb Q_\ell},\mathbb Z_\ell(w))$
when $w$ is odd, an isomorphism onto its torsion subgroup
when $\ell-1\mid w$, and zero in the remaining even weights.
The local group has free rank one in every case.

At $p\ne\ell$, evaluation at arithmetic Frobenius identifies
the local integral group with
$\mathbb Z_\ell/(p^w-1)\mathbb Z_\ell$.
In the middle case of E.210.1 the image of E.210.2 has exact
order $\ell^{1+v_\ell(w)}$ at every such $p$.
Thus localization of that global finite correction group
at any one off-diagonal prime is injective.

Consequently, for $w\ge3$, above a fixed class-$(w-1)$
section whose global lift fibre is nonempty, and with
actual local aligning frames retained, the matching
parameters are precisely the $b_w$-fold copies of
E.210.1. At an odd step, a prescribed compatible local
point section at $\ell$ selects a unique global lift class.
At a nonzero even step there is at most one lift matching
any prescribed nonempty set of support points away from
$\ell$; at the other even steps there is no global
correction parameter. These uniqueness assertions do not
assert that the other support matches hold.
\end{result}

\begin{proof}
The tame relation from A.20 kills the inertia value of
every integral weight-$w$ cocycle away from $\ell$:
$(p^w-p)c(T_p)=0$ and the coefficient module is torsion
free. The same holds for coboundaries. Inflation is
therefore an isomorphism, and A.20 makes the integral
group $H$ a compact finitely generated $\mathbb Z_\ell$-module.
The coefficient exact sequence and E.209 give
\[
 H/\ell H\simeq H^1(G_S,\mathbb F_\ell(w)).                    \tag{E.210.3}
\]
The global Euler characteristic
\citep[I.5.1, p.~67]{CohMilneADT}, with $H^2=0$, is
\[
 \dim H^0(G_S,\mathbb F_\ell(w))
 -\dim H^1(G_S,\mathbb F_\ell(w))
 =-\dim\mathbb F_\ell(w)^-.
\]
The right side is $-1$ in odd weight and zero in even
weight. The global invariant space has dimension one
exactly when $\ell-1\mid w$. Consequently the dimension
of E.210.3 is one in the first two cases of E.210.1
and zero in the third. Compact Nakayama gives a cyclic
module in the former cases and the zero module in the latter.

Since $w\ne0$, the rational Tate module has no invariants.
The usual invariant-coset map identifies the torsion of
$H$ with
\[
 (\mathbb Q_\ell(w)/\mathbb Z_\ell(w))^{G_S}.
                                                               \tag{E.210.4}
\]
Explicitly a fixed coset $v+\mathbb Z_\ell(w)$ maps to
the integral cocycle $g\mapsto g(v)-v$. Conversely, if
$\ell^a[c]=0$, then $\ell^ac=\delta t$ and
$c=\delta(t/\ell^a)$, proving that every torsion class
is obtained in this way. The global cyclotomic image
is $\mathbb Z_\ell^\times$. If $\ell-1\nmid w$, its
Teichm\"uller part contains $u$ with $u^w-1$ a unit,
so E.210.4 is zero. If $\ell-1\mid w$, the ideal
generated by all $u^w-1$ is
$\ell^{1+v_\ell(w)}\mathbb Z_\ell$: the binomial
expansion of $(1+\ell)^w$ gives exactly that valuation,
and all principal units give at least that valuation.
This proves the torsion order and E.210.2. A cyclic
module with this nonzero torsion is exactly this finite
cyclic group. In odd weight it is torsion-free and its
reduction is nonzero, so it is free of rank one.
This proves E.210.1 without a rational-rank substitution.

We check the integral local group and the comparison
map separately. For a nonzero integer $j$ put
\[
 a_j=\begin{cases}1+v_\ell(j),&\ell-1\mid j,\\0,&\ell-1\nmid j.
 \end{cases}
\]
The local cyclotomic image at $\ell$ is again
$\mathbb Z_\ell^\times$. Hence its finite invariants
in weight $j$ have order $\ell^{\min(n,a_j)}$.
Local duality and the local Euler characteristic
\citep[I.2.3 and I.2.8, pp.~28,31]{CohMilneADT} give
\[
 |H^1(G_{\mathbb Q_\ell},\mathbb Z/\ell^n(w))|
   =\ell^{n+\min(n,a_w)+\min(n,a_{1-w})}.
\]
Passing to the actual inverse limit in degree two,
with dual transitions the inclusions of torsion,
gives a group of order $\ell^{a_{1-w}}$. The local
coefficient exact sequence therefore gives
\[
 |H^1(G_{\mathbb Q_\ell},\mathbb Z_\ell(w))/\ell^n|
                  =\ell^{n+\min(n,a_w)}.                       \tag{E.210.5}
\]
All finite local groups are finite, so their derived
inverse-limit terms vanish. Compact Nakayama and
E.210.5 give free rank one. The same invariant-coset
argument identifies its torsion with a cyclic group
of order $\ell^{a_w}$. In particular the odd-weight
local integral group is free of rank one.

To prove that the odd-weight localization is a unit,
we use an actual finite injectivity statement.
Poitou--Tate duality
\citep[I.4.10(a), pp.~56--57]{CohMilneADT} identifies
the dual of the kernel of finite global localization
on $\mathbb F_\ell(w)$ with the kernel of localization
in degree two on its Tate dual
$\mathbb F_\ell(1-w)$. The latter group is zero by
E.209.5. At the real place all positive-degree
cohomology vanishes because $\ell$ is odd.
Thus
\[
 H^1(G_S,\mathbb F_\ell(w))
 \longrightarrow H^1(G_{\mathbb Q_\ell},\mathbb F_\ell(w))
 \quad\text{is injective}.                                    \tag{E.210.6}
\]
For odd $w$, E.210.3 and E.210.6 imply that the
localization of an integral global generator is
nonzero after reduction. The natural map from the
local integral group modulo $\ell$ into finite local
cohomology is injective by its coefficient exact
sequence. This is enough even when $w\equiv1\pmod{\ell-1}$
and finite local $H^1$ has dimension two rather than one.
The map between the integral free rank-one modules
is therefore multiplication by a unit, hence is an
isomorphism. In the nonzero even case E.210.2 comes
from the same invariant coset globally and locally,
and generates the whole local torsion subgroup.

At $p\ne\ell$ the integral tame value is zero by the
same torsion-free relation. Unramified cohomology
is the cokernel of $p^w-1$ on $\mathbb Z_\ell(w)$.
When $\ell-1\mid w$, its valuation is at least
$s_w=1+v_\ell(w)$, and the element
$(p^w-1)/\ell^{s_w}$ in that cokernel has exact
order $\ell^{s_w}$. This proves the stated injectivity.

Finally, for $w\ge3$, B.26 identifies the nonempty set
of actual lift classes above a fixed predecessor with a torsor under the
$b_w$ copies of global $H^1$. Its local aligning
frame is unique: a first nonzero weight-$i$ coordinate
of a geometric stabilizer would satisfy
$(\chi_\ell(g)^i-1)t_i=0$ for every local $g$,
where the local cyclotomic character has infinite
image and the coefficient module is torsion free.
This also holds at the coefficient place.
The localization descriptions now give exactly the
matching and uniqueness conclusions claimed.
No change of predecessor or unrecorded conjugator
is included in these assertions.
\end{proof}

At five the pattern is especially explicit: odd
weights have one $\mathbb Z_5$ parameter, weights
divisible by four have the finite parameter
$\mathbb Z/5^{1+v_5(w)}$, and weights congruent to
two modulo four have none. The particular sufficient
quintic condition of E.208 is bypassed by E.213's
two-first-prime construction. The infinite sequence
of support matches and the irregular coefficient
primes remain separate from this all-weight reduction.
